\documentclass[11pt]{article}
\usepackage[utf8]{inputenc}
\usepackage[T1]{fontenc}
\usepackage[english]{babel}
\usepackage{amsmath,amssymb}
\usepackage[a4paper,margin=1in]{geometry}
\usepackage{fancyhdr}
\usepackage{hyperref}
\usepackage[protrusion=true,expansion=false]{microtype}
\usepackage{array}

\pagestyle{fancy}
\fancyhf{}
\fancyhead[C]{\small Cayley, Collected Mathematical Papers --- Vol.~III, pp.~54--78 (Papers 173--177)}
\fancyfoot[C]{\thepage}
\renewcommand{\headrulewidth}{0.4pt}

\providecommand{\Sadlerian}{Sadlerian}
\providecommand{\degsign}{$^{\circ}$}
\providecommand{\dd}{d}

\hypersetup{
  pdftitle={Cayley Collected Mathematical Papers, Vol III, pp 54--78},
  pdfauthor={Arthur Cayley (transcription)},
  hidelinks
}

\setcounter{page}{54}

\allowdisplaybreaks

\begin{document}

%===============================================================
% PDF p.76 = book p.54 -- Paper 173 (cont.) "On Laplace's Method"
%===============================================================
\begin{center}
\small\textsc{54}\hfill\textsc{on laplace's method for the attraction of ellipsoids.}\hfill[173
\end{center}

\medskip

\noindent (and tending towards the centre) and for the potential of the entire ellipsoid the
well-known expressions
\[
A = 2 \int dS\, \frac{\xi\sqrt{R}}{L},
\]
\[
B = 2 \int dS\, \frac{\eta\sqrt{R}}{L},
\]
\[
C = 2 \int dS\, \frac{\zeta\sqrt{R}}{L},
\]
and
\[
V = 2 \int dS\, \frac{I\sqrt{R}\, dS}{L^{2}},
\]
where the integrations extend over the spherical angle of the circumscribed cone $R=0$.

We have moreover, by the general theory of attractions,
\[
A = -\frac{dV}{da},\quad B = -\frac{dV}{db},\quad C = -\frac{dV}{dc}.
\]

Since at the limits of the integration the quantity under the integral sign vanishes,
it is easy to see that the first differentials of $V$, $A$, $B$, $C$ with respect to any of the
quantities $a$, $b$, $c$, $l$, $m$, $n$, $k$, may be found by simply differentiating under the integral
sign without its being necessary to pay attention to the variation of the limits, so
that, for instance, $\dfrac{dV}{da} = 2 \int dS\, \dfrac{d}{da}\,\dfrac{I\sqrt{R}}{L}$. It is proper to remark that the expressions
thus obtained for the attractions $A$, $B$, $C$, are of a different form from the foregoing
expressions for the same quantities. Laplace writes
\[
F = aA + bB + cC,
\]
and he remarks that it may be shown by differentiation that the quantities $B$, $C$, $F$, $V$,
are connected by an equation which (writing $k$ for $k^{2}$, and $\dfrac{m}{l}$, $\dfrac{n}{l}$ for $m$, $n$, in the notation
of equation of the ellipsoid being with Laplace $x^{2} + my^{2} + nz^{2} = k^{2}$) becomes,
in the notation of the present memoir,
\[
\{l(a^{2} + b^{2} + c^{2}) - k\}\left( \frac{dV}{dk} - \frac{dF}{dk} \right) + V - F
\]
\[
+ \left(1 - \frac{l}{m}\right) b \left( \frac{dF}{db} - \tfrac{1}{2}\,\frac{dV}{db} - B\right)
\]
\[
+ \left(1 - \frac{l}{n}\right) c \left( \frac{dF}{dc} - \tfrac{1}{2}\,\frac{dV}{dc} - C\right)
\]
\[
- (m - l)\,\frac{dF}{dm} - (n - l)\,\frac{dF}{dn} = 0.
\]

\newpage

%===============================================================
% PDF p.77 = book p.55
%===============================================================
\begin{center}
\small 173]\hfill\textsc{on laplace's method for the attraction of ellipsoids.}\hfill\textsc{55}
\end{center}

\medskip

\noindent I write this equation in the form
\[
(la^{2} + mb^{2} + nc^{2} - k)\left( \frac{dV}{dk} - \frac{dF}{dk}\right) + V - F
\]
\[
+ (l - l)\left\{-a^{2}\left(\frac{dV}{dk} - \frac{dF}{dk}\right) + \frac{a}{l}\left(\frac{dF}{da} - \tfrac{1}{2}\,\frac{dV}{da} - A\right) - \frac{dF}{dl}\right\}
\]
\[
+ (m - l)\left\{-b^{2}\left(\frac{dV}{dk} - \frac{dF}{dk}\right) + \frac{b}{m}\left(\frac{dF}{db} - \tfrac{1}{2}\,\frac{dV}{db} - B\right) - \frac{dF}{dm}\right\}
\]
\[
+ (n - l)\left\{-c^{2}\left(\frac{dV}{dk} - \frac{dF}{dk}\right) + \frac{c}{n}\left(\frac{dF}{dc} - \tfrac{1}{2}\,\frac{dV}{dc} - C\right) - \frac{dF}{dn}\right\}
\]
\[
= 0;
\]
and I remark that this equation may be broken up into two equations, each of which
separately is satisfied; these equations are
\[
- k\left(\frac{dV}{dk} - \frac{dF}{dk}\right) + (V - F) - l\,\frac{dF}{dl} - m\,\frac{dF}{dm} - n\,\frac{dF}{dn}
\]
\[
+ a\left(\frac{dF}{da} - \tfrac{1}{2}\,\frac{dV}{da} - A\right) + b\left(\frac{dF}{db} - \tfrac{1}{2}\,\frac{dV}{db} - B\right) + c\left(\frac{dF}{dc} - \tfrac{1}{2}\,\frac{dV}{dc} - C\right) = 0,
\]
and
\[
(a^{2} + b^{2} + c^{2})\left(\frac{dV}{dk} - \frac{dF}{dk}\right) - \frac{dF}{dl} - \frac{dF}{dm} - \frac{dF}{dn}
\]
\[
+ \frac{a}{l}\left(\frac{dF}{da} - \tfrac{1}{2}\,\frac{dV}{da} - A\right) + \frac{b}{m}\left(\frac{dF}{db} - \tfrac{1}{2}\,\frac{dV}{db} - B\right) + \frac{c}{n}\left(\frac{dF}{dc} - \tfrac{1}{2}\,\frac{dV}{dc} - C\right) = 0.
\]

It may be added that the functions under the integral signs, and consequently
the integrals, are all of them homogeneous of the degree zero in $l$, $m$, $n$, $k$. The
thing to be verified is that the foregoing two equations are satisfied by the functions
under the integral signs, independently of the integrations, in fact by the values
\[
\left\{\begin{array}{l}
A = \dfrac{\xi\sqrt{R}}{L},\ B = \dfrac{\eta\sqrt{R}}{L},\ C = \dfrac{\zeta\sqrt{R}}{L}, \\[4pt]
V = \dfrac{I\sqrt{R}}{L}, \\[4pt]
F = \dfrac{(a\xi + b\eta + c\zeta)\sqrt{R}}{L}.
\end{array}\right.
\]

We find by differentiating these values, and after a few obvious reductions,
\[
\frac{dV}{da} = l\xi\left(\frac{\sqrt{R}}{L^{2}} + \frac{I}{L^{2}\sqrt{R}}\right) + I\sigma\left(\frac{-I}{L\sqrt{R}}\right),
\]
\[
\frac{dV}{dl} = a\xi\left(\frac{\sqrt{R}}{L^{2}} + \frac{I}{L^{2}\sqrt{R}}\right) + a^{2}\left(\frac{-I}{2L\sqrt{R}}\right) + \mathfrak{F}\left(-\frac{3I\sqrt{R}}{2L^{2}} - \frac{I^{3}}{2L^{2}\sqrt{R}}\right),
\]
\[
\frac{dV}{dk} = \frac{I}{2L\sqrt{R}},
\]

\newpage

%===============================================================
% PDF p.78 = book p.56
%===============================================================
\begin{center}
\small\textsc{56}\hfill\textsc{on laplace's method for the attraction of ellipsoids.}\hfill[173
\end{center}

\medskip

\[
\frac{dF}{da} \cdot \frac{\sqrt{R}}{L} = (a\xi + b\eta + c\zeta)\left\{l\xi\left(\frac{I}{L\sqrt{R}} + l\sigma\left(\frac{-1}{\sqrt{R}}\right)\right)\right\},
\]
\[
\frac{dF}{dl} = (a\xi + b\eta + c\zeta)\left\{a\xi\left(\frac{I}{L\sqrt{R}} + a^{2}\left(\frac{-1}{2\sqrt{R}}\right)\right) + \mathfrak{F}\left(-\frac{\sqrt{R}}{2L} - \frac{I^{2}}{2L\sqrt{R}}\right)\right\},
\]
\[
\frac{dF}{dk} = (a\xi + b\eta + c\zeta)\left\{\frac{1}{2L\sqrt{R}}\right\},
\]
values which give, as they should do,
\[
l\,\frac{dF}{dl} + m\,\frac{dF}{dm} + n\,\frac{dF}{dn} + k\,\frac{dF}{dk} = 0.
\]
Hence forming the values of the different parts of the first equation,
\[
- k\left(\frac{dV}{dk} - \frac{dF}{dk}\right) + V - F - l\,\frac{dF}{dl} - m\,\frac{dF}{dm} - n\,\frac{dF}{dn}
\]
\[
= \frac{I\sqrt{R}}{L^{2}} - \frac{kI}{2L\sqrt{R}} + (a\xi + b\eta + c\zeta)\left\{-\frac{\sqrt{R}}{L} + \frac{k}{L\sqrt{R}}\right\},
\]
\[
a\left(\frac{dF}{da} - A\right) + b\left(\frac{dF}{db} - B\right) + c\left(\frac{dF}{dc} - C\right) = (a\xi + b\eta + c\zeta)\left\{- \frac{\sqrt{R}}{L} + \frac{k}{L\sqrt{R}}\right\},
\]
\[
- \tfrac{1}{2}\left(a\,\frac{dV}{da} + b\,\frac{dV}{db} + c\,\frac{dV}{dc}\right) = - \frac{I\sqrt{R}}{L^{2}} + \frac{kI}{2L\sqrt{R}},
\]
which satisfy the first equation.

And in like manner for the second equation,
\[
- (a^{2} + b^{2} + c^{2})\left(\frac{dV}{dk} - \frac{dF}{dk}\right) = - (a^{2} + b^{2} + c^{2})\,\frac{I}{2L\sqrt{R}} + (a^{2} + b^{2} + c^{2})(a\xi + b\eta + c\zeta)\,\frac{1}{2L\sqrt{R}},
\]
\[
\frac{a}{l}\left(\frac{dF}{da} - A\right) + \frac{b}{m}\left(\frac{dF}{db} - B\right) + \frac{c}{n}\left(\frac{dF}{dc} - C\right)
\]
\[
= (a\xi + b\eta + c\zeta)\,\frac{I}{L\sqrt{R}} - (a^{2} + b^{2} + c^{2})(a\xi + b\eta + c\zeta)\,\frac{1}{2\sqrt{R}},
\]
\[
- \tfrac{1}{2}\left(\frac{a}{l}\,\frac{dV}{da} + \frac{b}{m}\,\frac{dV}{db} + \frac{c}{n}\,\frac{dV}{dc}\right) = - (a^{2} + b^{2} + c^{2})(a\xi + b\eta + c\zeta)\,\frac{I}{L\sqrt{R}} + (a\xi + b\eta + c\zeta)\left(\frac{\sqrt{R}}{2L^{2}} + \frac{I^{2}}{2L^{2}\sqrt{R}}\right),
\]
\[
- \frac{dF}{dl} - \frac{dF}{dm} - \frac{dF}{dn} = - (a\xi + b\eta + c\zeta)\,\frac{I}{2L\sqrt{R}} + (a^{2} + b^{2} + c^{2})(a\xi + b\eta + c\zeta)\,\frac{1}{2\sqrt{R}}
\]
\[
+ (a\xi + b\eta + c\zeta)\left(\frac{\sqrt{R}}{2L^{2}} + \frac{I^{2}}{2L^{2}\sqrt{R}}\right),
\]
which satisfy the second equation.

\newpage

%===============================================================
% PDF p.79 = book p.57
%===============================================================
\begin{center}
\small 173]\hfill\textsc{on laplace's method for the attraction of ellipsoids.}\hfill\textsc{57}
\end{center}

\medskip

\noindent Now considering $V$, $F$, $A$, $B$, $C$ as standing for the definite integrals, we may
replace $A$, $B$, $C$ by the differential coefficients of $V$, and retaining $F$ for shortness to
stand for
\[
F = - a\,\frac{dV}{da} - b\,\frac{dV}{db} - c\,\frac{dV}{dc}\;;
\]
the first equation becomes
\[
- k\left(\frac{dV}{dk} - \frac{dF}{dk}\right) + V - F - l\,\frac{dF}{dl} - m\,\frac{dF}{dm} - n\,\frac{dF}{dn}
\]
\[
+ a\left(\frac{dF}{da} + \tfrac{1}{2}\,\frac{dV}{da}\right) + b\left(\frac{dF}{db} + \tfrac{1}{2}\,\frac{dV}{db}\right) + c\left(\frac{dF}{dc} + \tfrac{1}{2}\,\frac{dV}{dc}\right) = 0;
\]
and the second equation becomes
\[
- (a^{2} + b^{2} + c^{2})\left(\frac{dV}{dk} - \frac{dF}{dk}\right) - \frac{dF}{dl} - \frac{dF}{dm} - \frac{dF}{dn}
\]
\[
+ \frac{a}{l}\left(\frac{dF}{da} + \tfrac{1}{2}\,\frac{dV}{da}\right) + \frac{b}{m}\left(\frac{dF}{db} + \tfrac{1}{2}\,\frac{dV}{db}\right) + \frac{c}{n}\left(\frac{dF}{dc} + \tfrac{1}{2}\,\frac{dV}{dc}\right) = 0.
\]

If we put as before $V = \dfrac{I\sqrt{R}}{L}$, the preceding values of the differential coefficients
of $V$ give $F = \dfrac{2I\sqrt{R}}{L^{2}} + \dfrac{kI}{L\sqrt{R}}$, or as we may write it $F = -2V + W$, where $W = \dfrac{kI}{L\sqrt{R}}$.

I put then for the moment
\[
\left\{
\begin{array}{l}
V = \dfrac{I\sqrt{R}}{L},\quad W = \dfrac{kI}{L\sqrt{R}}, \\[6pt]
F = - 2V + W.
\end{array}
\right.
\]

It should be remarked that there is nothing in what has preceded which tends
to show that these values must satisfy the differential equations. The definite integrals
must, of course, satisfy as before the equations, but it does not follow that the equations
are satisfied by the elements separately. And in fact only the first equation is so
satisfied; the second equation is not satisfied. To verify this I form the differential
coefficients
\[
\frac{dW}{da} = l\xi\left(\frac{k}{L\sqrt{R}} - \frac{kl}{LR\sqrt{R}}\right) + l\sigma\left(\frac{kl}{R\sqrt{R}}\right),
\]
\[
\frac{dW}{dk} = \frac{I}{L\sqrt{R}} - \frac{kl}{2LR\sqrt{R}},
\]
\[
\frac{dW}{dl} = a\xi\left(\frac{k}{L\sqrt{R}} - \frac{kl}{LR\sqrt{R}}\right) + \mathfrak{F}\left(\frac{-3kl}{2L^{2}\sqrt{R}} + \frac{kI^{2}}{2L^{2}R\sqrt{R}}\right) + a^{2}\,\frac{kl}{2R\sqrt{R}}.
\]

\bigskip
\hfill\textsc{c. iii.}\hfill 8

\newpage

%===============================================================
% PDF p.80 = book p.58
%===============================================================
\begin{center}
\small\textsc{58}\hfill\textsc{on laplace's method for the attraction of ellipsoids.}\hfill[173
\end{center}

\medskip

\noindent We have as before
\[
l\,\frac{dF}{dl} + m\,\frac{dF}{dm} + n\,\frac{dF}{dn} + k\,\frac{dF}{dk} = 0,
\]
and forming the expressions for the different parts of the first equation,
\[
- k\left(\frac{dV}{dk} - \frac{dF}{dk}\right) + V - F - l\,\frac{dF}{dl} - m\,\frac{dF}{dm} - n\,\frac{dF}{dn}
\]
\[
= - 5k\,\frac{dV}{dk} + 5k\,\frac{dW}{dk} + 3V - W = -\frac{3I\sqrt{R}}{L^{2}} - \frac{3kl}{2L\sqrt{R}} - \frac{kI}{L\sqrt{R}},
\]
\[
a\left(\frac{dF}{da} + \tfrac{1}{2}\,\frac{dV}{da}\right) + \&c.\, =
\]
\[
- \tfrac{3}{2}\left(a\,\frac{dV}{da} + b\,\frac{dV}{db} + c\,\frac{dV}{dc}\right) + \left(a\,\frac{dW}{da} + b\,\frac{dW}{db} + c\,\frac{dW}{dc}\right) = -\frac{3I\sqrt{R}}{L^{2}} + \frac{3kl}{2L\sqrt{R}},
\]
and
\[
- \tfrac{1}{2}\left(a\,\frac{dV}{da} + b\,\frac{dV}{db} + c\,\frac{dV}{dc}\right) = -\frac{kl}{L\sqrt{R}},
\]
which show that the first equation is satisfied.

Proceeding in the same manner with the second equation,
\[
- (a^{2} + b^{2} + c^{2})\left(\frac{dV}{dk} - \frac{dF}{dk}\right) =
\]
\[
- (a^{2} + b^{2} + c^{2})\left(3\,\frac{dV}{dk} - \frac{dW}{dk}\right) = (a^{2} + b^{2} + c^{2})\left(-\frac{I}{L\sqrt{R}} - \frac{kl}{2RL\sqrt{R}}\right),
\]
\[
\frac{a}{l}\left(\frac{dF}{da} - A\right) + \frac{b}{m}\left(\frac{dF}{db} - B\right) + \frac{c}{n}\left(\frac{dF}{dc} - C\right) =
\]
\[
- \tfrac{1}{2}\left(\frac{a}{l}\,\frac{dV}{da} + \frac{b}{m}\,\frac{dV}{db} + \frac{c}{n}\,\frac{dV}{dc}\right) + \left(\frac{a}{l}\,\frac{dW}{da} + \frac{b}{m}\,\frac{dW}{db} + \frac{c}{n}\,\frac{dW}{dc}\right) = (a^{2} + b^{2} + c^{2})\left(- \frac{3\sqrt{R}}{2L^{2}} - \frac{3I^{2}}{2L^{2}\sqrt{R}}\right);
\]
\[
- \frac{dF}{dl} - \frac{dF}{dm} - \frac{dF}{dn} = 2\left(\frac{dV}{dl} + \frac{dV}{dm} + \frac{dV}{dn}\right) - \left(\frac{dW}{dl} + \frac{dW}{dm} + \frac{dW}{dn}\right),
\]

\newpage

%===============================================================
% PDF p.81 = book p.59
%===============================================================
\begin{center}
\small 173]\hfill\textsc{on laplace's method for the attraction of ellipsoids.}\hfill\textsc{59}
\end{center}

\medskip

\noindent and
\[
2\left(\frac{dV}{dl} + \frac{dV}{dm} + \frac{dV}{dn}\right) = (a^{2} + b^{2} + c^{2})\left(\frac{-1}{L\sqrt{R}}\right) + (a\xi + b\eta + c\zeta)\left(\frac{2\sqrt{R}}{L^{2}} + \frac{2I^{2}}{L^{2}\sqrt{R}}\right) - \frac{3I\sqrt{R}}{L^{2}} - \frac{I^{2}}{L^{2}\sqrt{R}},
\]
\[
- \left(\frac{dW}{dl} + \frac{dW}{dm} + \frac{dW}{dn}\right) = (a^{2} + b^{2} + c^{2})\left(\frac{-kI}{2RL\sqrt{R}}\right) + (a\xi + b\eta + c\zeta)\left(\frac{-k}{L\sqrt{R}} + \frac{kI^{2}}{LR\sqrt{R}}\right)
\]
\[
+ \frac{3kl}{2L\sqrt{R}} + \frac{kI^{2}}{2L\sqrt{R}\,R};
\]
and the value of the left-hand side of the second equation is therefore,
\[
(a\xi + b\eta + c\zeta)\left(\frac{\sqrt{R}}{2L^{2}} + \frac{3I^{2}}{2L^{2}\sqrt{R}}\right) - \frac{3I\sqrt{R}}{L^{2}} - \frac{I^{2}}{2L^{2}\sqrt{R}} + \frac{3kl}{2L\sqrt{R}} - \frac{kl^{2}}{2L\sqrt{R}\,R},
\]
which is not equal to zero, or the second equation is not satisfied.

I consider again $V$, $F$ as denoting the definite integrals, and I eliminate $\dfrac{dV}{dl}$, $\dfrac{dF}{dl}$
by means of the equations
\[
l\,\frac{dV}{dl} + m\,\frac{dV}{dm} + n\,\frac{dV}{dn} + k\,\frac{dV}{dk} = 0,
\]
\[
l\,\frac{dF}{dl} + m\,\frac{dF}{dm} + n\,\frac{dF}{dn} + k\,\frac{dF}{dk} = 0.
\]

The first equation thus becomes
\[
l\,\frac{dV}{dl} = \frac{dV}{dm} + n\,\frac{dV}{dn} - 2\left(l\,\frac{dF}{dl} + m\,\frac{dF}{dm} + n\,\frac{dF}{dn}\right) + V - F
\]
\[
+ \tfrac{1}{2}\left(a\,\frac{dV}{da} + b\,\frac{dV}{db} + c\,\frac{dV}{dc}\right) + \left(a\,\frac{dF}{da} + b\,\frac{dF}{db} + c\,\frac{dF}{dc}\right) = 0,
\]
and the second equation becomes
\[
(a^{2} + b^{2} + c^{2})\,\frac{1}{l}\left(l\,\frac{dV}{dl} + m\,\frac{dV}{dm} + n\,\frac{dV}{dn}\right) - (a^{2} + b^{2} + c^{2})\,\frac{1}{l}\left(l\,\frac{dF}{dl} + m\,\frac{dF}{dm} + n\,\frac{dF}{dn}\right)
\]
\[
+ \tfrac{1}{2}\left(\frac{a}{l}\,\frac{dV}{da} + \frac{b}{m}\,\frac{dV}{db} + \frac{c}{n}\,\frac{dV}{dc}\right) + \left(\frac{a}{l}\,\frac{dF}{da} + \frac{b}{m}\,\frac{dF}{db} + \frac{c}{n}\,\frac{dF}{dc}\right) - \left(\frac{dF}{dl} + \frac{dF}{dm} + \frac{dF}{dn}\right) = 0\;;
\]
and it will be remembered that in these equations
\[
F = - a\,\frac{dV}{da} - b\,\frac{dV}{db} - c\,\frac{dV}{dc}.
\]

The first equation (it is easy to perceive) shows merely that $V$ is made up of terms
separately homogeneous in $a$, $b$, $c$, and in $l$, $m$, $n$, and such that the degrees in the
two sets respectively being $\kappa$, $\lambda$, then $\lambda = \tfrac{1}{2}(\kappa - 2)$. In fact $V$ being a function of
the form in question, if we attend only to the term the degrees of which are $\kappa$, $\tfrac{1}{2}\lambda$,

\medskip
\hfill 8---2

\newpage

%===============================================================
% PDF p.82 = book p.60
%===============================================================
\begin{center}
\small\textsc{60}\hfill\textsc{on laplace's method for the attraction of ellipsoids.}\hfill[173
\end{center}

\medskip

\noindent then, by the properties of homogeneous functions, $F = - V$, and the first equation is
satisfied if only
\[
\lambda + 2\lambda\kappa + 1 + \kappa + \tfrac{1}{2}\kappa - \kappa^{2} = 0,
\]
i.e.\ if $\lambda(2\kappa+1) = \kappa^{2} - \tfrac{3}{2}\kappa - 1 = \tfrac{1}{2}(\kappa - 2)(2\kappa + 1)$ or $\lambda = \tfrac{1}{2}(\kappa - 2)$. Or, what is the same thing,
we may say that the first equation shows that $V$ is made up of terms the degrees
of which in $a$, $b$, $c$ and in $l$, $m$, $n$, are $\kappa$, $-2i-1$, and $i$, $-i-\tfrac{1}{2}$ respectively. Attending
henceforth only to the second equation, I write $\dfrac{l}{k} = \dfrac{1}{\alpha + \theta}$, $\dfrac{m}{k} = \dfrac{1}{\beta + \theta}$, $\dfrac{n}{k} = \dfrac{1}{\gamma + \theta}$; so
that $\alpha + \theta$, $\beta + \theta$, $\gamma + \theta$ denote the squared semiaxes of the ellipsoid. We have
\[
\frac{d}{dl} = -\frac{(\alpha + \theta)^{2}}{k}\,\frac{d}{d\alpha},\quad \&c.,
\]
and the equation becomes
\[
- (a^{2} + b^{2} + c^{2})\left((\alpha + \theta)\,\frac{dV}{d\alpha} + (\beta + \theta)\,\frac{dV}{d\beta} + (\gamma + \theta)\,\frac{dV}{d\gamma}\right)
\]
\[
+ (a^{2} + b^{2} + c^{2})\left((\alpha + \theta)\,\frac{dF}{d\alpha} + (\beta + \theta)\,\frac{dF}{d\beta} + (\gamma + \theta)\,\frac{dF}{d\gamma}\right)
\]
\[
+ \tfrac{1}{2}\left(a(\alpha + \theta)\,\frac{dV}{da} + b(\beta + \theta)\,\frac{dV}{db} + c(\gamma + \theta)\,\frac{dV}{dc}\right)
\]
\[
+ \left(a(\alpha + \theta)\,\frac{dF}{da} + b(\beta + \theta)\,\frac{dF}{db} + c(\gamma + \theta)\,\frac{dF}{dc}\right)
\]
\[
+ \left((\alpha + \theta)^{2}\,\frac{dF}{d\alpha} + (\beta + \theta)^{2}\,\frac{dF}{d\beta} + (\gamma + \theta)^{2}\,\frac{dF}{d\gamma}\right) = 0.
\]

Put for shortness $\Theta = (\alpha + \theta)(\beta + \theta)(\gamma + \theta)$, and write
\[
V = v\sqrt{\Theta},\quad F = f\sqrt{\Theta},
\]
($\sqrt{\Theta}$ is to a constant factor \emph{pr\`es} the mass of the ellipsoid) then $v$, $f$ are connected
by the equation
\[
f = - a\,\frac{dv}{da} - b\,\frac{dv}{db} - c\,\frac{dv}{dc};
\]
and observing that
\[
(\alpha + \theta)\,\frac{d\sqrt{\Theta}}{d\alpha} + (\beta + \theta)\,\frac{d\sqrt{\Theta}}{d\beta} + (\gamma + \theta)\,\frac{d\sqrt{\Theta}}{d\gamma} = 3\Theta\,\frac{1}{2\sqrt{\Theta}} = \tfrac{3}{2}\sqrt{\Theta},
\]
\[
(\alpha + \theta)^{2}\,\frac{d\sqrt{\Theta}}{d\alpha} + (\beta + \theta)^{2}\,\frac{d\sqrt{\Theta}}{d\beta} + (\gamma + \theta)^{2}\,\frac{d\sqrt{\Theta}}{d\gamma}
\]
\[
= (\alpha + \beta + \gamma + 3\theta)\,\Theta\,\frac{1}{2\sqrt{\Theta}} = \tfrac{1}{2}(\alpha + \beta + \gamma + 3\theta)\sqrt{\Theta},
\]

\newpage

%===============================================================
% PDF p.83 = book p.61
%===============================================================
\begin{center}
\small 173]\hfill\textsc{on laplace's method for the attraction of ellipsoids.}\hfill\textsc{61}
\end{center}

\medskip

\noindent the equation becomes
\[
- (a^{2} + b^{2} + c^{2})\left((\alpha + \theta)\,\frac{dv}{d\alpha} + (\beta + \theta)\,\frac{dv}{d\beta} + (\gamma + \theta)\,\frac{dv}{d\gamma}\right)
\]
\[
+ (a^{2} + b^{2} + c^{2})\left((\alpha + \theta)\,\frac{df}{d\alpha} + (\beta + \theta)\,\frac{df}{d\beta} + (\gamma + \theta)\,\frac{df}{d\gamma}\right)
\]
\[
- \tfrac{1}{2}(a^{2} + b^{2} + c^{2})\,v + \tfrac{1}{2}(a^{2} + b^{2} + c^{2})\,f
\]
\[
+ \tfrac{1}{2}\left(a(\alpha + \theta)\,\frac{dv}{da} + b(\beta + \theta)\,\frac{dv}{db} + c(\gamma + \theta)\,\frac{dv}{dc}\right)
\]
\[
+ \left(a(\alpha + \theta)\,\frac{df}{da} + b(\beta + \theta)\,\frac{df}{db} + c(\gamma + \theta)\,\frac{df}{dc}\right)
\]
\[
+ \left((\alpha + \theta)^{2}\,\frac{df}{d\alpha} + (\beta + \theta)^{2}\,\frac{df}{d\beta} + (\gamma + \theta)^{2}\,\frac{df}{d\gamma}\right)
\]
\[
+ \tfrac{1}{2}(\alpha + \beta + \gamma + 3\theta)\,f = 0.
\]

Now $v\left(=\dfrac{1}{\sqrt{\Theta}}\,V\right)$ may be expanded in the form
\[
v = U_{0} + U_{1} + U_{2} + \cdots,
\]
where $U_{i}$ is of the degree $2i$ in the semiaxes $\sqrt{(\alpha + \theta)}$, $\sqrt{(\beta + \theta)}$, $\sqrt{(\gamma + \theta)}$, and of the
degree $-2i-1$ in the coordinates $a$, $b$, $c$. And it is easy to see that the first term
of the expansion is
\[
U_{0} = \frac{4\pi}{3}\cdot\frac{1}{\sqrt{(a^{2} + b^{2} + c^{2})}}.
\]

The preceding value of $v$ gives
\[
f = U_{0} + 3U_{1} + 5U_{2} + \cdots + (2i + 1)\,U_{i} + \cdots,
\]
and substituting the values of $v$, $f$ in the differential equation, and attending only to
the terms which are of the same dimensions as $(a^{2} + b^{2} + c^{2})\,U_{i+1}$, we have
\[
\left\{- (i + 1) + (i + 1)(2i + 3) - \tfrac{1}{2} + \tfrac{1}{2}(2i + 3)\right\}(a^{2} + b^{2} + c^{2})\,U_{i+1}
\]
\[
+ \left\{\alpha x\,\frac{dU_{i}}{d\alpha} + b\beta\,\frac{dU_{i}}{d\beta} + c\gamma\,\frac{dU_{i}}{d\gamma} - (2i + 1)\,\theta U_{i}\right\}\left\{\tfrac{1}{2} + 2i + 1\right\}
\]
\[
+ \left\{a^{2}\,\frac{dU_{i}}{d\alpha} + \beta^{2}\,\frac{dU_{i}}{d\beta} + \gamma^{2}\,\frac{dU_{i}}{d\gamma} + 2i\theta U_{i} - \theta^{2}\left(\frac{dU_{i}}{d\alpha} + \frac{dU_{i}}{d\beta} + \frac{dU_{i}}{d\gamma}\right)\right\}(2i + 1)
\]
\[
+ \tfrac{1}{2}(\alpha + \beta + \gamma + 3\theta)(2i + 1)\,U_{i} = 0;
\]

\newpage

%===============================================================
% PDF p.84 = book p.62
%===============================================================
\begin{center}
\small\textsc{62}\hfill\textsc{on laplace's method for the attraction of ellipsoids.}\hfill[173
\end{center}

\medskip

\noindent or, reducing,
\[
(i + 1)(2i + 5)(a^{2} + b^{2} + c^{2})\,U_{i+1}
\]
\[
+ (2i + \tfrac{1}{2})\left(\alpha a\,\frac{dU_{i}}{da} + b\beta\,\frac{dU_{i}}{db} + c\gamma\,\frac{dU_{i}}{dc}\right)
\]
\[
+ (2i + 1)\left(\alpha^{2}\,\frac{dU_{i}}{d\alpha} + \beta^{2}\,\frac{dU_{i}}{d\beta} + \gamma^{2}\,\frac{dU_{i}}{d\gamma} + \tfrac{1}{2}(\alpha + \beta + \gamma)\,U_{i}\right)
\]
\[
- (2i + 1)\left(\frac{dU_{i}}{d\alpha} + \frac{dU_{i}}{d\beta} + \frac{dU_{i}}{d\gamma}\right)\theta^{2} = 0.
\]

Now $U_{i}$ is a function of $\alpha + \theta$, $\beta + \theta$, $\gamma + \theta$; if, therefore, for any particular value
of $i$, $U_{i}$ is independent of $\theta$, it is clear that $U_{i}$ must be a function of the differences
of these quantities, and consequently we shall have $\dfrac{dU_{i}}{d\alpha} + \dfrac{dU_{i}}{d\beta} + \dfrac{dU_{i}}{d\gamma} = 0$; and this being so, the
equation of differences, and consequently $U_{i+1}$, will be independent of $\theta$. But $U_{0}$ is
independent of $\theta$, hence $U_{1}$, $U_{2}$, $\&c.$ are all of them independent of $\theta$, or $v$ is in\-dependent of $\theta$; i.e.\ for ellipsoids having the same foci for their principal sections, and
acting on the same external point, the potentials, and therefore the attractions, are
proportional to the masses, which is Maclaurin's theorem for the attractions of ellipsoids
upon an external point.

The foregoing equation, omitting the term which vanishes, gives
\[
U_{i+1} = -\,\frac{(2i + 1)\left\{\alpha a\,\dfrac{dU_{i}}{da} + b\beta\,\dfrac{dU_{i}}{db} + c\gamma\,\dfrac{dU_{i}}{dc}\right\} + (2i + 1)\left(\alpha^{2}\,\dfrac{dU_{i}}{d\alpha} + \beta^{2}\,\dfrac{dU_{i}}{d\beta} + \gamma^{2}\,\dfrac{dU_{i}}{d\gamma} + \tfrac{1}{2}(\alpha + \beta + \gamma)\,U_{i}\right)}{(i + 1)(2i + 5)(a^{2} + b^{2} + c^{2})}.
\]

It may be remarked that this equation, with the assistance of the equation
\[
\frac{dU_{i}}{d\alpha} + \frac{dU_{i}}{d\beta} + \frac{dU_{i}}{d\gamma} = 0,
\]
gives
\[
(i + 1)(2i + 5)\left(\frac{dU_{i+1}}{d\alpha} + \frac{dU_{i+1}}{d\beta} + \frac{dU_{i+1}}{d\gamma}\right) = - (2i + \tfrac{1}{2})\left(a\,\frac{dU_{i}}{da} + b\,\frac{dU_{i}}{db} + c\,\frac{dU_{i}}{dc}\right)
\]
\[
= (2i + 1)\left(2\left(\alpha\,\frac{dU_{i}}{d\alpha} + \beta\,\frac{dU_{i}}{d\beta} + \gamma\,\frac{dU_{i}}{d\gamma}\right) + \tfrac{3}{2}U_{i}\right) = \left((2i + \tfrac{1}{2})(2i + 1) - (2i + 1)(2i + \tfrac{1}{2})\right)U_{i} = 0,
\]
which is as it should be.

Write
\[
U_{i} = \frac{Q_{i}}{(a^{2} + b^{2} + c^{2})^{i + \tfrac{1}{2}}}.
\]

\newpage

%===============================================================
% PDF p.85 = book p.63
%===============================================================
\begin{center}
\small 173]\hfill\textsc{on laplace's method for the attraction of ellipsoids.}\hfill\textsc{63}
\end{center}

\medskip

\noindent where $U_{i}$ is of the degree $i$ in $\alpha$, $\beta$, $\gamma$, and of the degree $2i$ in $a$, $b$, $c$. We have
\[
a\alpha\,\frac{dU_{i}}{d\alpha} + b\beta\,\frac{dU_{i}}{d\beta} + c\gamma\,\frac{dU_{i}}{dc} = (a^{2} + b^{2} + c^{2})^{-i - \tfrac{1}{2}}\left(\alpha a\,\frac{dQ_{i}}{d\alpha} + b\beta\,\frac{dQ_{i}}{d\beta} + c\gamma\,\frac{dQ_{i}}{dc}\right)
\]
\[
- (4i + 1)(a^{2} + b^{2} + c^{2})^{-i - \tfrac{3}{2}}(a^{2}\alpha + b^{2}\beta + c^{2}\gamma)\,Q_{i},
\]
\[
\alpha^{2}\,\frac{dU_{i}}{d\alpha} + \beta^{2}\,\frac{dU_{i}}{d\beta} + \gamma^{2}\,\frac{dU_{i}}{d\gamma} + \tfrac{1}{2}(\alpha + \beta + \gamma)\,U_{i}
\]
\[
= (a^{2} + b^{2} + c^{2})^{-i - \tfrac{1}{2}}\left(\alpha^{2}\,\frac{dQ_{i}}{d\alpha} + \beta^{2}\,\frac{dQ_{i}}{d\beta} + \gamma^{2}\,\frac{dQ_{i}}{d\gamma}\right) + \tfrac{1}{2}(\alpha + \beta + \gamma)\,Q_{i}.
\]

Hence, putting in like manner
\[
(i + 1)(2i + 5)\,Q_{i+1} = - (2i + \tfrac{1}{2})\left\{(a^{2} + b^{2} + c^{2})\left(\alpha a\,\frac{dQ_{i}}{d\alpha} + b\beta\,\frac{dQ_{i}}{d\beta} + c\gamma\,\frac{dQ_{i}}{dc}\right)\right.
\]
\[
\left. - (4i + 1)(a^{2}\alpha + b^{2}\beta + c^{2}\gamma)\,Q_{i}\right\}
\]
\[
- (2i + 1)(a^{2} + b^{2} + c^{2})\left(\alpha^{2}\,\frac{dQ_{i}}{d\alpha} + \beta^{2}\,\frac{dQ_{i}}{d\beta} + \gamma^{2}\,\frac{dQ_{i}}{d\gamma}\right) + \tfrac{1}{2}(\alpha + \beta + \gamma)\,Q_{i},
\]
i.e.
\[
(i + 1)(2i + 5)\,Q_{i+1} = - (a^{2} + b^{2} + c^{2})\left\{(2i + \tfrac{1}{2})\left(\alpha a\,\frac{dQ_{i}}{d\alpha} + b\beta\,\frac{dQ_{i}}{d\beta} + c\gamma\,\frac{dQ_{i}}{dc}\right)\right.
\]
\[
\left. + (2i + 1)\left(\alpha^{2}\,\frac{dQ_{i}}{d\alpha} + \beta^{2}\,\frac{dQ_{i}}{d\beta} + \gamma^{2}\,\frac{dQ_{i}}{d\gamma} + \tfrac{1}{2}(\alpha + \beta + \gamma)\,Q_{i}\right)\right\} + (2i + 1)(4i + 1)(a^{2}\alpha + b^{2}\beta + c^{2}\gamma)\,Q_{i},
\]
from which the functions $Q_{i}$ may be calculated successively.

We may, it is clear, write
\[
U_{i} = \frac{4\pi}{3}\cdot\frac{1}{2^{i}\cdot 1.2.3 \cdots i\cdot 5.7 \cdots 2i + 3}\,K_{i},
\]
and we shall then have
\[
(a^{2} + b^{2} + c^{2})\,K_{i+1} + (4i + 3)\left(\alpha a\,\frac{dK_{i}}{d\alpha} + b\beta\,\frac{dK_{i}}{db} + c\gamma\,\frac{dK_{i}}{dc}\right)
\]
\[
+ (4i + 2)\left(\alpha^{2}\,\frac{dK_{i}}{d\alpha} + \beta^{2}\,\frac{dK_{i}}{d\beta} + \gamma^{2}\,\frac{dK_{i}}{d\gamma}\right) + (2i + 1)(\alpha + \beta + \gamma)\,K_{i} = 0,
\]
where
\[
K_{0} = \frac{1}{\sqrt{(a^{2} + b^{2} + c^{2})}}.
\]
I proceed to show that
\[
K_{i} = \left(\alpha\,\frac{d^{2}}{da^{2}} + \beta\,\frac{d^{2}}{db^{2}} + \gamma\,\frac{d^{2}}{dc^{2}}\right)^{i}\frac{1}{\sqrt{(a^{2} + b^{2} + c^{2})}}.
\]

\newpage

%===============================================================
% PDF p.86 = book p.64
%===============================================================
\begin{center}
\small\textsc{64}\hfill\textsc{on laplace's method for the attraction of ellipsoids.}\hfill[173
\end{center}

\medskip

\noindent In fact, assuming this equation for any particular value of $i$, we find first
\[
\left(\alpha^{2}\,\frac{d}{d\alpha} + \beta^{2}\,\frac{d}{d\beta} + \gamma^{2}\,\frac{d}{d\gamma}\right)K_{i} = i\left(\alpha\,\frac{d^{2}}{da^{2}} + \beta\,\frac{d^{2}}{db^{2}} + \gamma\,\frac{d^{2}}{dc^{2}}\right)^{i-1}\left(\alpha^{2}\,\frac{d^{2}}{da^{2}} + \beta^{2}\,\frac{d^{2}}{db^{2}} + \gamma^{2}\,\frac{d^{2}}{dc^{2}}\right)\frac{1}{\sqrt{(a^{2} + b^{2} + c^{2})}}
\]
\[
= i\left(\alpha\,\frac{d^{2}}{da^{2}} + \beta\,\frac{d^{2}}{db^{2}} + \gamma\,\frac{d^{2}}{dc^{2}}\right)^{i-1}\left(\alpha\,\frac{d^{2}}{da^{2}} + \beta\,\frac{d^{2}}{db^{2}} + \gamma\,\frac{d^{2}}{dc^{2}}\right)\frac{1}{\sqrt{(a^{2} + b^{2} + c^{2})}}
\]
\[
= i\left(\alpha\,\frac{d^{2}}{da^{2}} + \beta\,\frac{d^{2}}{db^{2}} + \gamma\,\frac{d^{2}}{dc^{2}}\right)^{i}K_{i-1}.
\]

Now putting for shortness
\[
\Delta = \alpha\,\frac{d^{2}}{da^{2}} + \beta\,\frac{d^{2}}{db^{2}} + \gamma\,\frac{d^{2}}{dc^{2}},
\]
we find, replacing $\Delta K_{i+1}$ and $\Delta K_{i}$ by their values $K_{i+2}$ and $K_{i+1}$,
\[
\Delta\left((a^{2} + b^{2} + c^{2})\,K_{i+1}\right) = (a^{2} + b^{2} + c^{2})\,K_{i+2}
\]
\[
+ 4\left(\alpha a\,\frac{d}{da} + b\beta\,\frac{d}{db} + c\gamma\,\frac{d}{dc}\right)K_{i+1}
\]
\[
+ 2\,(\alpha + \beta + \gamma)\,K_{i+1},
\]
\[
\Delta\left(\alpha a\,\frac{d}{da} + b\beta\,\frac{d}{db} + c\gamma\,\frac{d}{dc}\right)K_{i} = \left(\alpha a\,\frac{d}{da} + b\beta\,\frac{d}{db} + c\gamma\,\frac{d}{dc}\right)K_{i+1}
\]
\[
+ 2\left(\alpha^{2}\,\frac{d^{2}}{da^{2}} + \beta^{2}\,\frac{d^{2}}{db^{2}} + \gamma^{2}\,\frac{d^{2}}{dc^{2}}\right)K_{i},
\]
\[
\Delta\left(\alpha^{2}\,\frac{d}{d\alpha} + \beta^{2}\,\frac{d}{d\beta} + \gamma^{2}\,\frac{d}{d\gamma}\right)K_{i} = \left(\alpha^{2}\,\frac{d}{d\alpha} + \beta^{2}\,\frac{d}{d\beta} + \gamma^{2}\,\frac{d}{d\gamma}\right)K_{i+1}
\]
\[
- \left(\alpha^{2}\,\frac{d^{2}}{da^{2}} + \beta^{2}\,\frac{d^{2}}{db^{2}} + \gamma^{2}\,\frac{d^{2}}{dc^{2}}\right)K_{i},
\]
\[
\Delta(\alpha + \beta + \gamma)\,K_{i} = (\alpha + \beta + \gamma)\,K_{i+1};
\]
hence operating on the equation of differences with the symbol $\Delta$, we obtain
\[
(a^{2} + b^{2} + c^{2})\,K_{i+2} + (4i + 7)\left(\alpha a\,\frac{d}{da} + b\beta\,\frac{d}{db} + c\gamma\,\frac{d}{dc}\right)K_{i+1}
\]
\[
+ (4i + 2)\left(\alpha^{2}\,\frac{d}{d\alpha} + \beta^{2}\,\frac{d}{d\beta} + \gamma^{2}\,\frac{d}{d\gamma}\right)K_{i+1}
\]
\[
+ \{2(4i + 3) - (4i + 2)\}\left(\alpha^{2}\,\frac{d^{2}}{da^{2}} + \beta^{2}\,\frac{d^{2}}{db^{2}} + \gamma^{2}\,\frac{d^{2}}{dc^{2}}\right)K_{i}
\]
\[
+ (2i + 3)\,(\alpha + \beta + \gamma)\,K_{i+1} = 0,
\]

\newpage

%===============================================================
% PDF p.87 = book p.65
%===============================================================
\begin{center}
\small 173]\hfill\textsc{on laplace's method for the attraction of ellipsoids.}\hfill\textsc{65}
\end{center}

\medskip

\noindent the third line of which is
\[
4(i + 1)\left(\alpha^{2}\,\frac{d^{2}}{da^{2}} + \beta^{2}\,\frac{d^{2}}{db^{2}} + \gamma^{2}\,\frac{d^{2}}{dc^{2}}\right)K_{i} = 4i\left(\alpha^{2}\,\frac{d}{d\alpha} + \beta^{2}\,\frac{d}{d\beta} + \gamma^{2}\,\frac{d}{d\gamma}\right)K_{i+1},
\]
by a foregoing equation, and the assumed equation of difference thus leads to
\[
(a^{2} + b^{2} + c^{2})\,K_{i+2} + (4i + 7)\left(\alpha a\,\frac{d}{da} + b\beta\,\frac{d}{db} + c\gamma\,\frac{d}{dc}\right)K_{i+1}
\]
\[
+ (4i + 6)\left(\alpha^{2}\,\frac{d}{d\alpha} + \beta^{2}\,\frac{d}{d\beta} + \gamma^{2}\,\frac{d}{d\gamma}\right)K_{i+1}
\]
\[
+ (2i + 3)\,(\alpha + \beta + \gamma)\,K_{i+1} = 0,
\]
which is the assumed equation, writing $i + 1$ instead of $i$. The equation, if true for $i$,
is therefore true for $i + 1$, and it is easily seen to be true for $i = 0$; hence it is true
generally, or the value
\[
K_{i} = \left(\alpha\,\frac{d^{2}}{da^{2}} + \beta\,\frac{d^{2}}{db^{2}} + \gamma\,\frac{d^{2}}{dc^{2}}\right)^{i}\,\frac{1}{\sqrt{(a^{2} + b^{2} + c^{2})}}
\]
satisfies the equation obtained by Laplace's method, and gives, moreover, the proper
value for $K_{0}$. We have thus the value of $K_{i}$; and remembering that
\[
V = \sqrt{(\alpha + \theta)(\beta + \theta)(\gamma + \theta)}\,v,
\]
and observing that the symbol $\Delta$ may be replaced by
\[
\Delta = (\alpha + \theta)\,\frac{d^{2}}{da^{2}} + (\beta + \theta)\,\frac{d^{2}}{db^{2}} + (\gamma + \theta)\,\frac{d^{2}}{dc^{2}},
\]
the value of $V$ is
\[
V = \frac{4\pi}{3}\sqrt{(\alpha + \theta)(\beta + \theta)(\gamma + \theta)}\,\sum_{0}^{\infty}\!\left(\frac{1}{2^{i}\cdot 1\cdot 2 \cdots i\cdot 5\cdot 7 \cdots 2i + 3}\,\Delta^{i}\,\frac{1}{\sqrt{(a^{2} + b^{2} + c^{2})}}\right),
\]
which is in fact the value which I have found by a much more simple method in
the \emph{Cambridge Mathematical Journal}, t.~\textsc{iii}.\ p.~69 [2].

\bigskip
\hfill\textsc{c. iii.}\hfill 9

\newpage

%===============================================================
% PDF p.88 = book p.66 -- Paper 174 "On the Oval of Descartes"
%===============================================================
\begin{center}
\small\textsc{66}\hfill\hfill[174
\end{center}

\vspace{4em}

\begin{center}
{\large\textbf{174.}}\\[1em]
{\large\scshape On the Oval of Descartes.}\\[1em]
{\small [From the \textit{Quarterly Mathematical Journal}, vol.~\textsc{i}.\ (1857), pp.~320--328.]}
\end{center}

\bigskip

\textsc{This} is an extract of the passage in the Geometry (\textit{La G\'eom\'etrie de M.\ Descartes}, 12\textsuperscript{mo}\ Paris, 1705,
pp.~79--92) in which he applies his method of coordinates to this now well-known curve: the marginal
reference is ``\textit{Explication de quatre nouveaux genres d'Ovales qui servent \`a l'Optique}.'' The paper was intended
to be introductory to a discussion in some detail of the history and properties of the Curve, but no continuation was written.

\newpage

%===============================================================
% PDF p.89 = book p.67 -- Paper 175 "Porism of the In-and-Circumscribed Triangle"
%===============================================================
\begin{center}
\small 175]\hfill\hfill\textsc{67}
\end{center}

\vspace{4em}

\begin{center}
{\large\textbf{175.}}\\[1em]
{\large\scshape On the Porism of the In-and-Circumscribed Triangle.}\\[1em]
{\small [From the \textit{Quarterly Mathematical Journal}, vol.~\textsc{i}.\ (1857), pp.~344--354.]}
\end{center}

\bigskip

\textsc{The} porism of the in-and-circumscribed triangle in its most general form relates
to a triangle the angles of which lie in fixed curves, and the sides of which touch
fixed curves; but at present I consider only the case in which the angles lie in one
and the same fixed curve, which for greater simplicity I assume to be a conic. We
have therefore a triangle $ABC$, the angles of which lie in a fixed conic $\mathfrak{S}$, and the
sides of which touch the fixed curves $\mathfrak{A}$, $\mathfrak{B}$, $\mathfrak{C}$; the points of contact may be repre\-sented by $\alpha$, $\beta$, $\gamma$. And if we consider the conic $\mathfrak{S}$ and the curves $\mathfrak{A}$, $\mathfrak{B}$ as given,
the curve $\mathfrak{C}$ will be the envelope of the side $AB$; to construct this side we have
only to take at pleasure a point $C$ on the curve $\mathfrak{S}$ and to draw through this point
tangents to the curves $\mathfrak{B}$, $\mathfrak{A}$ respectively meeting the conic $\mathfrak{S}$ in the points $A$ and $B$;
the line joining these points is the required side $AB$. I may notice that in the
case supposed of the curve $\mathfrak{S}$ being a conic, the lines $A\alpha$, $B\beta$, $C\gamma$ meet in a point---
which gives at once a construction for $\gamma$, the point of contact of $AB$ with the curve $\mathfrak{C}$.
For the sake however of exhibiting the reasoning in a form which may be modified
so as to be applicable to a curve $\mathfrak{S}$ of any order, instead of the conic $\mathfrak{S}$, I shall
dispense with the employment of the property just mentioned, which is peculiar to
the case of the conic.

Suppose for a moment (figs.~1 and 1 \emph{bis}) that the curves $\mathfrak{A}$, $\mathfrak{B}$ are points, and let
the line through $M$, meet $\mathfrak{A}$, $\mathfrak{B}$ the conic $\mathfrak{S}$ in the points $M$, $N$. If we take the point
$N$ for the angle $C$ of the triangle, the points $A$, $B$ will each of them coincide with
$M$, and the side $AB$ will be the tangent at $M$ to the conic $\mathfrak{S}$: call this tangent $T$.
Consider next a point $C$ in the neighbourhood of $N$, we shall have two points $A$, $B$
in the neighbourhood of $M$, and the point in which $T$ intersects $AB$ will be the
point of contact $\gamma$ of $T$ with the curve $\mathfrak{C}$. To find this point, suppose that $M\mathfrak{A} = a$,
$N\mathfrak{A} = a'$, $M\mathfrak{B} = b$, $N\mathfrak{B} = b'$ and let the distance of $C$ from $N$ be $ds$; the distances
of $A$, $B$ from the line $MN$ parallel to $T$ will be proportional to $\dfrac{b}{b'}\,ds$, $\dfrac{a}{a'}\,ds$, and

\medskip
\hfill 9---2

\newpage

%===============================================================
% PDF p.90 = book p.68
%===============================================================
\begin{center}
\small\textsc{68}\hfill\textsc{on the porism of the in-and-circumscribed triangle.}\hfill[175
\end{center}

\medskip

\noindent the perpendicular distances of these points from $T$ are consequently proportional to
$\dfrac{b^{2}}{b'^{2}}\,ds^{2}$, $\dfrac{a^{2}}{a'^{2}}\,ds^{2}$. The inclination of $AB$ to $T$ is therefore proportional to
\[
\left(\frac{b^{2}}{b'^{2}} - \frac{a^{2}}{a'^{2}}\right)ds^{2} = \left(\frac{b}{b'} - \frac{a}{a'}\right)ds,\ \text{i.e.\ to}\ \left(\frac{b}{b'} + \frac{a}{a'}\right)ds,
\]
which is of the same order as $ds$, and it is at once seen that the point $\gamma$ will
coincide with $M$, i.e.\ that the curve $\mathfrak{C}$ touches the conic $\mathfrak{S}$ at the point $M$. If,
however, $\dfrac{a}{a'} + \dfrac{b}{b'} = 0$, i.e.\ if the points $M$, $N$ are harmonically related to $\mathfrak{A}$, $\mathfrak{B}$, then the
inclination of $AB$ to $T$ is in general of the order $ds^{2}$, and the point $\gamma$ will be at a finite distance
from $M$; moreover $T$ is in this case a stationary tangent (i.e.\ a tangent at a point
of inflexion) of the curve $\mathfrak{C}$. Now reverting to the general case of $\mathfrak{A}$, $\mathfrak{B}$ being any
two curves, then if there be a common tangent touching these curves in the points
$\alpha$, $\beta$, and meeting the curve $\mathfrak{S}$ in the points $M$, $N$, the like reasonings apply to this
case. Hence

\textsc{First Lemma.} If a common tangent to the curves $\mathfrak{A}$, $\mathfrak{B}$ touch these curves in
$\alpha$, $\beta$ and meet the conic $\mathfrak{S}$ in the points $M$, $N$; the point $N$ gives rise to a branch
of the curve $\mathfrak{C}$ which (except in the case after mentioned) touches the conic $\mathfrak{S}$ at
the point $M$. If however $M$, $N$ are harmonically related with respect to $\alpha$, $\beta$, then
the branch of the curve $\mathfrak{C}$ does not pass through $M$, but it has for a stationary
tangent the tangent $M$ to the conic $\mathfrak{S}$.

Suppose again that the curve $\mathfrak{A}$ (fig.~2) is a point, and let the curve $\mathfrak{B}$ intersect
the conic $\mathfrak{S}$ at the point $M$, and let the tangent to $\mathfrak{B}$ at $M$ meet $\mathfrak{S}$ in a point $R$,
and $R\mathfrak{A}$ meet $\mathfrak{S}$ in a point $Q$. Then taking the point $R$ for the angle $C$ of the
triangle, we shall have $A$, $B$ coinciding with $M$, $Q$ respectively, and thence $MQ$ a
tangent to the curve $\mathfrak{C}$. To find the point of contact, take $C$ in the neighbourhood
of $R$ at a distance from it $ds$; $B$ will be a point in the neighbourhood of $Q$ and
distant from it by an infinitesimal of the order $ds$, $A$ will be in the neighbourhood
of $M$ and distant from it by an infinitesimal of the order $ds^{2}$. Hence $AB$ will
intersect $MQ$ at the point $M$, or the curve $\mathfrak{C}$ will pass through $M$. Reverting to the
general case where $\mathfrak{A}$ is a curve, we have only to consider $RQ$ as a tangent to the
curve $\mathfrak{A}$ at a point $\alpha$, and the like reasoning will apply to this case: hence

\newpage

%===============================================================
% PDF p.91 = book p.69
%===============================================================
\begin{center}
\small 175]\hfill\textsc{on the porism of the in-and-circumscribed triangle.}\hfill\textsc{69}
\end{center}

\medskip

\textsc{Second Lemma.} If the curve $\mathfrak{B}$ intersect the conic $\mathfrak{S}$ at a point $M$, and the
tangent to $\mathfrak{B}$ at $M$ intersect $\mathfrak{S}$ in $R$, if moreover a tangent through $R$ to the
curve $\mathfrak{A}$ intersect $\mathfrak{S}$ in $Q$; then to each of the points $Q$ there corresponds a branch
through $M$ of the curve $\mathfrak{C}$, viz.\ a branch touching $MQ$ at the point $M$.

Suppose as before (fig.~3) that $\mathfrak{A}$ is a point, and let the curve $\mathfrak{B}$ intersect the
conic $\mathfrak{S}$ at the point $M$, and the tangent to $\mathfrak{B}$ at $M$ meet $\mathfrak{S}$ in the point $R$. And let
$M\mathfrak{A}$ meet $\mathfrak{S}$ in the point $N$. Then taking the point $R$ for the angle $C$ of the
triangle we shall have $A$, $B$ coinciding with $R$, $N$ respectively, and thence $NR$ a tangent
to the curve $\mathfrak{C}$. To find the point of contact, take $C$ in the neighbourhood of $R$
at a distance from it $ds$, then $C$ will be in the neighbourhood of $M$ and distant
from it by an infinitesimal of the order $ds^{2}$, and $B$ will be in the neighbourhood
of $N$ and distant from it by an infinitesimal of the same order $ds^{2}$, and consequently
$AB$ will intersect $NR$ at the point $N$, or the curve $\mathfrak{C}$ passes through $N$. Reverting
to the general case where $\mathfrak{A}$ is a curve, we have only to consider $MN$ as a tangent
to the curve $\mathfrak{A}$ at a point $\alpha$ and the like reasoning applies: hence,

\textsc{Third Lemma.} If the curve $\mathfrak{B}$ intersect the conic $\mathfrak{S}$ at a point $M$, and the
tangent to $\mathfrak{B}$ at $M$ intersect $\mathfrak{S}$ in $R$, if moreover a tangent through $M$ to the curve

\newpage

%===============================================================
% PDF p.92 = book p.70
%===============================================================
\begin{center}
\small\textsc{70}\hfill\textsc{on the porism of the in-and-circumscribed triangle.}\hfill[175
\end{center}

\medskip

\noindent $\mathfrak{A}$ meet $\mathfrak{S}$ in $N$; then to each of the points $R$ there corresponds a branch through
$N$ of the curve $\mathfrak{C}$, viz.\ a branch touching $NR$ at $N$.

Suppose again (fig.~4) that the curve $\mathfrak{A}$ is a point, and let the curve $\mathfrak{B}$ touch
$\mathfrak{S}$ at the point $M$. And let $M\mathfrak{A}$ meet $\mathfrak{S}$ in the point $N$. Then taking the point
$M$ for the angle $C$ of the triangle, the angle $A$ will coincide with $M$ and the angle
$B$ will coincide with $N$, and consequently we shall have $MN$ a tangent to the curve
$\mathfrak{C}$. And $MN$ is in fact a double tangent, for proceeding to find the point of contact,
take $C$ in the neighbourhood of $M$ at a distance $ds$; then from $C$ we may draw to
the curve $\mathfrak{B}$ two tangents each of them meeting $\mathfrak{S}$ in a point $\mathfrak{A}$ in the neighbour\-hood of $M$ and distant from it by an infinitesimal of the order $ds$; again $CA$ meets
$\mathfrak{S}$ in a point $B$ in the neighbourhood of $N$ and distant from it by an infinitesimal
of the same order $ds$; hence $AB$ will meet $MN$ at a point which will be in general
at a finite distance from $M$, or rather (since there are two positions of the point $A$)
the lines $AB$ will meet $MN$ in two points, each of them in general at a finite
distance from $M$. Reverting to the case of $\mathfrak{A}$ being a curve, we must as before
consider $MN$ as a tangent to the curve $\mathfrak{A}$ at the point $\alpha$, and the same reasoning
applies: hence

\textsc{Fourth Lemma.} If the curve $\mathfrak{B}$ touch the conic $\mathfrak{S}$ at the point $M$, and if a
tangent through $M$ to the curve $\mathfrak{A}$ meet $\mathfrak{S}$ in $N$, then $MN$ is a double tangent of
the curve $\mathfrak{C}$, viz.\ the line $MN$ has two distinct points of contact with the curve $\mathfrak{C}$.

Suppose (fig.~5) that the curves $\mathfrak{A}$ and $\mathfrak{B}$ meet the conic $\mathfrak{S}$ in one and the
same point $M$, and let the tangent at $M$ to the curve $\mathfrak{A}$ meet $\mathfrak{S}$ in the point $P$,
and the tangent at $M$ to the curve $\mathfrak{B}$ meet $\mathfrak{S}$ in the point $N$; then if we take
the point $M$ for the angle $C$ of the triangle, the angles $A$ and $B$ of the triangle
will coincide with $N$, $P$ respectively, and $NP$ will be a tangent of the curve $\mathfrak{C}$.
But $NP$ will be a double tangent, for proceeding to find the point of contact, take
$C$ in the neighbourhood of $M$, and distant from it by an infinitesimal of the second
order $ds^{2}$; then since from the point $C$ there may be drawn two tangents touching

\newpage

%===============================================================
% PDF p.93 = book p.71
%===============================================================
\begin{center}
\small 175]\hfill\textsc{on the porism of the in-and-circumscribed triangle.}\hfill\textsc{71}
\end{center}

\medskip

\noindent $\mathfrak{A}$ in the neighbourhood of $M$, and two tangents touching $\mathfrak{B}$ in the neighbourhood of
$M$, there will be two points $A$ in the neighbourhood of $N$ and distant from it by
infinitesimals of the order $ds$, and in like manner two points $B$ in the neighbourhood
of $P$ and distant from it by infinitesimals of the order $ds$. Call these points $A$, $A'$;
$B$, $B'$; then $AB$, $A'B'$ will meet $NP$ in one and the same point, and $AB'$, $A'B$
will in like manner meet $NP$ in one and the same point; these two points, which
will be in general at finite distances from $N$, $P$, will be points of contact of $NP$
with the curve $\mathfrak{C}$: hence,

\textsc{Fifth Lemma.} If the curves $\mathfrak{A}$, $\mathfrak{B}$ meet the conic $\mathfrak{S}$ in one and the same
point $M$, and if the tangents at $M$ to the curves $\mathfrak{B}$ and $\mathfrak{A}$ respectively meet $\mathfrak{S}$ in
the points $N$ and $P$, then, joining these points, the line $NP$ will be a double
tangent to the curve $\mathfrak{C}$, viz.\ there will be two distinct points of contact of the line
$NP$ with the curve $\mathfrak{C}$.

The double tangents of the fourth and fifth lemmas exist by virtue of the par\-ticular relations assumed between the curves $\mathfrak{A}$, $\mathfrak{B}$ and the conic $\mathfrak{S}$, viz.\ from one
or both of the curves $\mathfrak{A}$, $\mathfrak{B}$ touching the conic $\mathfrak{S}$, or from the two curves having
a common point of intersection or common points of intersection with $\mathfrak{S}$; there are
besides double tangents which exist independently of any such relations, and the theory
of which will be presently investigated, but it will be convenient in the first instance
to find the class of $\mathfrak{C}$.

\newpage

%===============================================================
% PDF p.94 = book p.72
%===============================================================
\begin{center}
\small\textsc{72}\hfill\textsc{on the porism of the in-and-circumscribed triangle.}\hfill[175
\end{center}

\medskip

\noindent The class of $\mathfrak{C}$ is at once determinable from the classes (which I represent by
$m$, $n$) of the curves $\mathfrak{A}$, $\mathfrak{B}$. In fact, take any point $M$ on the conic $\mathfrak{S}$ for the
angle $A$ of the triangle. Through the point $M$ we may draw $n$ tangents to $\mathfrak{B}$,
each of which will intersect the conic $\mathfrak{S}$ in a single point, or there will be $n$
positions of the point $C$, and since from each of these we may draw $m$ tangents to
the curve $\mathfrak{A}$, each tangent intersecting $\mathfrak{S}$ in a single point; or we have $mn$ positions
of the angle $B$, i.e.\ this same number of tangents through $M$ to the curve $\mathfrak{C}$. But if
the same point $M$ had been taken as the angle $B$ of the triangle, there would have
been in like manner $mn$ positions of the angle $A$, i.e.\ this same number of tangents
through $M$ to the curve $\mathfrak{C}$. Hence there are in all $2mn$ tangents through $M$ to the
curve $\mathfrak{C}$, or the curve $\mathfrak{C}$ is of the class $2mn$.

Considering now the double tangents of the curve $\mathfrak{C}$; imagine a quadrilateral
inscribed in the conic $\mathfrak{S}$ of which two opposite sides touch the curve $\mathfrak{A}$, and the
other two opposite sides touch the curve $\mathfrak{B}$: suppose $MNPQ$, of which the sides $MN$
and $PQ$ touch $\mathfrak{A}$ and the sides $NP$ and $QM$ touch $\mathfrak{B}$. If we take $M$ for the angle
$C$ of the triangle, then the angles $A$, $B$ will coincide with $Q$, $N$, i.e.\ $NQ$ is a tangent
to the curve $\mathfrak{C}$, and the point of contact may be determined as before by considering
a point $C$ in the neighbourhood of $M$; but in like manner if we take $P$ for the
angle $C$ of the triangle, then the angles $A$, $B$ will coincide with $N$, $Q$, i.e.\ $QN$ is
again a tangent to the curve $\mathfrak{C}$, and the point of contact may be determined by
considering a point $C$ in the neighbourhood of $P$; $QN$ is therefore a double tangent
of the curve $\mathfrak{C}$. But in like manner $MP$ is a double tangent of the curve $\mathfrak{C}$, i.e.\
the diagonals of the quadrilateral are each of them double tangents of the curve $\mathfrak{C}$,
and the number of double tangents is consequently double the number of quadrilaterals.

Imagine a pentilateral inscribed in the conic $\mathfrak{S}$, the first and third sides of which
touch the curve $\mathfrak{A}$ and the second and fourth sides of which touch the curve $\mathfrak{B}$,
the fifth or closing side will envelope a curve $\mathfrak{C}'$, and in the cases in which the
curve $\mathfrak{C}'$ and the conic $\mathfrak{S}$ have a common tangent, the fifth side will vanish, and
the pentilateral becomes a quadrilateral of the kind before referred to. Now as $\mathfrak{C}$
was shown to be of the class $2mn$, so it may be shown that $\mathfrak{C}'$ is of the class
$4mn - m - n$, hence $\mathfrak{C}'$ and $\mathfrak{S}$ have $4m^{2}n^{2} - 2mn$ common tangents. The quadrilateral may reduce
itself to a common tangent of the curves $\mathfrak{A}$ and $\mathfrak{B}$: this gives rise to $2mn$ points
of contact of $\mathfrak{C}$ and $\mathfrak{S}$, and the common tangent at a point of contact reckons as
two common tangents; the number of the remaining common tangents is therefore
$4m^{2}n^{2} - 4mn$. And these are in fact tangents at points of contact of $\mathfrak{C}'$ and $\mathfrak{S}$;
i.e.\ $\mathfrak{C}'$ and $\mathfrak{S}$ touch in $2m^{2}n^{2} - 2mn$ points. And since each angle of the quadrilateral
may be taken as the first angle, the number of quadrilaterals is one fourth of this,
or $\tfrac{1}{2}(m^{2}n^{2} - mn)$, and the number of double tangents of the curve $\mathfrak{C}$ from the before-
mentioned cause is therefore $m^{2}n^{2} - mn$.

But there is another way in which double tangents arise; we may have a quad\-rilateral $MNPQ$ inscribed in the conic $\mathfrak{S}$, such that two adjacent sides $MN$, $NP$
touch the curve $\mathfrak{A}$, and the other two adjacent sides $PQ$, $QM$ touch the curve $\mathfrak{B}$.
In fact in this case one of the diagonals, viz.\ $NQ$, is a double tangent of the curve
$\mathfrak{C}$; the number of double tangents is therefore equal to the number of quadrilaterals.

\newpage

%===============================================================
% PDF p.95 = book p.73
%===============================================================
\begin{center}
\small 175]\hfill\textsc{on the porism of the in-and-circumscribed triangle.}\hfill\textsc{73}
\end{center}

\medskip

\noindent Consider a pentilateral inscribed in the conic, and such that the first and second
sides touch the curve $\mathfrak{A}$ and the third and fourth sides touch the curve $\mathfrak{B}$, the fifth
or closing side will envelope a curve $\mathfrak{C}''$, and in the cases in which the curve $\mathfrak{C}''$
and the conic $\mathfrak{S}$ have a common tangent, the fifth side will vanish, and the penti\-lateral will become a quadrilateral of the kind last referred to. The curve $\mathfrak{C}''$ is of
the class $2mn - (m - 1)(n - 1)$; hence $\mathfrak{C}''$ and $\mathfrak{S}$ have $4mn - (m - 1)(n - 1)$ common tangents;
but these are tangents at points of contact, or the curves touch in $2mn - (m - 1)(n - 1)$
points, and the number of quadrilaterals is $mn(m - 1)(n - 1)$. The number of double
tangents of the curve $\mathfrak{C}$ from the cause last referred to is therefore $mn(m - 1)(n - 1)$,
which is equal to $mn(mn - m - n + 1)$. The total number of double tangents of the
curve $\mathfrak{C}$ is consequently $mn(2mn - m - n)$. And the curve $\mathfrak{C}$ has not in general any
stationary tangents or what is the same thing any inflexions. It has been shown that
$\mathfrak{C}$ is of the class $2mn$, it is therefore of the order $2mn(2mn - 1) - 2mn(2mn - m - n)$,
which is equal to $2mn(m + n - 1)$. Hence

\textsc{Theorem.} If a triangle $ABC$ be inscribed in a conic $\mathfrak{S}$, and the sides $BC$, $AC$
touch curves $\mathfrak{A}$, $\mathfrak{B}$ of the classes $m$, $n$ respectively, the side $AB$ will envelope a
curve $\mathfrak{C}$ of the class $2mn$ with in general $mn(2mn - m - n)$ double tangents, but no
stationary tangents, and therefore of the order $2mn(m + n - 1)$. If the curve $\mathfrak{A}$ touch
the conic $\mathfrak{S}$, each point of contact will give rise to $n$ double tangents of the curve
$\mathfrak{C}$, and so if the curve $\mathfrak{B}$ touch the conic $\mathfrak{S}$, each point of contact will give rise to
$m$ double tangents of the curve $\mathfrak{C}$. And if $\mathfrak{A}$ and $\mathfrak{B}$ intersect on the conic $\mathfrak{S}$, each
such intersection will give rise to a double tangent of the curve $\mathfrak{C}$. The curve $\mathfrak{C}$
in general touches the conic $\mathfrak{S}$ in the points in which it is intersected by any
common tangent of the curves $\mathfrak{A}$ and $\mathfrak{B}$; but if the points of contact be harmonically
situated with respect to the conic $\mathfrak{S}$, then $\mathfrak{C}$ does not pass through the points of
intersection, but the tangents to $\mathfrak{S}$ at the points of intersection are stationary tangents
of $\mathfrak{C}$. There is of course in the above-mentioned special cases a corresponding re\-duction in the order of $\mathfrak{C}$.

It sometimes happens that the number of double tangents of $\mathfrak{C}$ becomes infinite,
i.e.\ in fact that $\mathfrak{C}$ is made up of two coincident curves; instances of this will be
presently mentioned.

Suppose that the curves $\mathfrak{A}$, $\mathfrak{B}$ are points; $\mathfrak{C}$ is in general a conic having double
contact with the conic $\mathfrak{S}$ in the points in which it is intersected by the line joining
$\mathfrak{A}$, $\mathfrak{B}$. But if the points $\mathfrak{A}$, $\mathfrak{B}$ are harmonically situated with respect to the conic $\mathfrak{S}$,
then $\mathfrak{C}$ does not pass through the points of intersection, but the tangents to $\mathfrak{S}$ at
the points of intersection are stationary tangents of $\mathfrak{C}$. This implies that the curve
$\mathfrak{C}$ is made up of two coincident points at the point of intersection of the two
tangents of $\mathfrak{S}$: call this the point $\mathfrak{C}$, then $\mathfrak{A}$, $\mathfrak{B}$, $\mathfrak{C}$ are conjugate points of the
conic $\mathfrak{S}$, and we have the well-known theorem that in a conic $\mathfrak{S}$ there may be
inscribed an infinity of triangles the sides of which pass through three conjugate
points of the conic. It should be remarked that for each position of the side $AB$,
there are two positions of the angle $C$, i.e.\ each side $AB$ is properly a double tan\-gent of the curve (point) $\mathfrak{C}$.

\bigskip
\hfill\textsc{c. iii.}\hfill 10

\newpage

%===============================================================
% PDF p.96 = book p.74
%===============================================================
\begin{center}
\small\textsc{74}\hfill\textsc{on the porism of the in-and-circumscribed triangle.}\hfill[175
\end{center}

\medskip

\noindent Let $\mathfrak{A}$ be a point and $\mathfrak{B}$ be a conic, the curve $\mathfrak{C}$ is in general of the class $4$,
with two double tangents, and therefore of the order $8$. It is proper to remark that
the double tangents originate in a quadrilateral of the kind first considered, viz.\ a
quadrilateral of which two opposite sides pass through the point $\mathfrak{A}$ and the other
two opposite sides touch the conic $\mathfrak{B}$. It is easy in the present case to construct
the quadrilateral: consider the point $\mathfrak{A}$ as a pole, and take its polar with respect to
the conic $\mathfrak{S}$, take the pole of this polar with respect to the conic $\mathfrak{B}$. Join the two
poles, and from the point in which the joining line meets the polar draw tangents
to the conic $\mathfrak{B}$, these tangents will meet the conic $\mathfrak{S}$ in four points, lying two and
two in lines passing through the point $\mathfrak{A}$: we have thus the quadrilateral, and the
diagonals of the quadrilateral (or the other two lines passing through the four points)
are the double tangents of the curve $\mathfrak{C}$.

There are two particular cases to be considered. First, where the conic $\mathfrak{B}$ has
double contact with the conic $\mathfrak{S}$. Here the lines joining the point $\mathfrak{A}$ with the points
of contact are double tangents of the curve $\mathfrak{C}$, which has therefore in all four double
tangents, and being a curve of the fourth class it must break up into two curves
of the second class, i.e.\ into two conics. And these are curves having double contact
with the conic $\mathfrak{S}$; for the curve $\mathfrak{C}$ touches the conic $\mathfrak{S}$ in the four points in which
$\mathfrak{S}$ is intersected by the tangents through $\mathfrak{A}$ to the conic $\mathfrak{B}$. Secondly, where $\mathfrak{A}$ is
one of the conjugate points of the conics $\mathfrak{B}$, $\mathfrak{S}$. The general construction for the
quadrilateral shows that if from any point of the common polar of $\mathfrak{A}$ with respect to
$\mathfrak{B}$ and with respect to $\mathfrak{S}$, we draw tangents to $\mathfrak{B}$, these will meet $\mathfrak{S}$ in four points,
lying two and two in lines passing through $\mathfrak{A}$, i.e.\ that the number of the double
tangents of the curve $\mathfrak{C}$ is indefinite; $\mathfrak{C}$ is therefore made up of two coincident
curves of the second class, i.e.\ of two coincident conics. Moreover, $\mathfrak{C}$ passes through
the points of intersection of $\mathfrak{B}$, $\mathfrak{S}$; hence, disregarding one of the two coincident conics,
we may say that the curve $\mathfrak{C}$ is a conic passing through the points of intersection of
the conics $\mathfrak{B}$, $\mathfrak{S}$.

Next, let the curves $\mathfrak{A}$, $\mathfrak{B}$ be each of them a conic. The curve $\mathfrak{C}$ is of the
class $8$, with in general $16$ double tangents, and therefore of the order $24$. But
there are two particular cases to be considered: first, where the conics $\mathfrak{A}$, $\mathfrak{B}$ have
each of them double contact with the conic $\mathfrak{S}$. Here the tangents drawn from the
points of contact of either of the conics $\mathfrak{A}$ or $\mathfrak{B}$ with $\mathfrak{S}$ to the other of the conics,
$\mathfrak{B}$ or $\mathfrak{A}$, is a double tangent of the curve $\mathfrak{C}$, i.e.\ there are $8$ new double tangents,
or in all $24$ double tangents of the curve $\mathfrak{C}$, which is therefore of the order $8$;
and being of the class $8$ with $24$ double tangents, it must break up into four curves
of the second class, i.e.\ into four conics. And the curve $\mathfrak{C}$ touches the conic $\mathfrak{S}$ in
the points in which $\mathfrak{S}$ is intersected by any one of the four common tangents of
the conics $\mathfrak{A}$, $\mathfrak{B}$, viz.\ $8$ points in all; hence each of the four conics has double
contact with the conic $\mathfrak{S}$. Attending only to one of the four conics of which $\mathfrak{C}$ is
made up, we have thus what (in a restricted sense of the expression, the porism of
the in-and-circumscribed triangle) I call the porism (homographic) of the in-and-
circumscribed triangle, viz.

\newpage

%===============================================================
% PDF p.97 = book p.75
%===============================================================
\begin{center}
\small 175]\hfill\textsc{on the porism of the in-and-circumscribed triangle.}\hfill\textsc{75}
\end{center}

\medskip

\noindent If a triangle be inscribed in a conic, and two of the sides touch conics having
double contact with the circumscribed conic, then will the third side touch a conic
having double contact with the circumscribed conic.

Secondly, the conics $\mathfrak{A}$ and $\mathfrak{B}$ may cut the conic $\mathfrak{S}$ in the same four points.
Here it may be seen that there are an infinity of inscribed quadrilaterals of the
kind first considered, viz.\ of which two opposite sides touch the conic $\mathfrak{A}$, and the
other two opposite sides touch the conic $\mathfrak{B}$. Hence, the curve $\mathfrak{C}$ is made up of
two coincident curves of the class $4$. But the curve of the class $4$ has in fact $4$
double tangents, viz.\ considering each of the points of intersection of $\mathfrak{A}$, $\mathfrak{B}$, $\mathfrak{S}$, and
drawing tangents to $\mathfrak{A}$ and $\mathfrak{B}$ meeting $\mathfrak{S}$ in two new points, the line joining these
points is a double tangent of the curve in question, which is therefore of the 4th
order, and being of the class $4$ with $4$ double tangents, it must break up into two
curves of the second class, i.e.\ into two conics. Each of these conics passes through
the points of intersection of $\mathfrak{A}$, $\mathfrak{B}$, $\mathfrak{S}$, and touches the four lines last referred to,
the conics would of course be completely determined by the condition of passing
through the four points and touching one of the four lines. Attending only to one
of the two conics, we have thus what I call the porism (allographic) of the in-and-
circumscribed triangle, viz.

If a triangle be inscribed in a conic, and two of the sides touch conics meeting
the circumscribed conic in the same four points, the remaining side will touch a
conic meeting the circumscribed conic in the four points.

The \emph{a posteriori} demonstration of these theorems will form the subject of an\-other paper, [178].

\bigskip
\hfill 10---2

\newpage

%===============================================================
% PDF p.98 = book p.76 -- Paper 176 "Note on Jacobi's Canonical Formulae"
%===============================================================
\begin{center}
\small\textsc{76}\hfill\hfill[176
\end{center}

\vspace{4em}

\begin{center}
{\large\textbf{176.}}\\[1em]
{\large\scshape Note on Jacobi's Canonical Formul\ae{} for Disturbed\\ Motion in an Elliptic Orbit.}\\[1em]
{\small [From the \textit{Quarterly Mathematical Journal}, vol.~\textsc{i}.\ (1857), pp.~355--356.]}
\end{center}

\bigskip

\textsc{Consider} a body (afterwards called the disturbed body) revolving about a central
body under the influence of their mutual attraction and of any disturbing forces.
Then referring the disturbed body to axes through the central body and parallel to
fixed lines in space, write
\begin{align*}
& x,\ y,\ z,\quad\text{the coordinates of the disturbed body,}\\
& r, \quad\text{the radius vector, } = \sqrt{(x^{2} + y^{2} + z^{2})},\\
& M, \quad\text{the mass of the disturbed body,}\\
& M', \quad\text{the mass of the central body.}
\end{align*}
Write also

$R$, the disturbing function, taken negatively, i.e.\ the sign of $R$ is taken as in
the \emph{M\'ecanique C\'eleste}.

The equations of motion then are
\[
\frac{d^{2}x}{dt^{2}} = - \frac{(M + M')\,x}{r^{3}} - \frac{dR}{dx},
\]
\[
\frac{d^{2}y}{dt^{2}} = - \frac{(M + M')\,y}{r^{3}} - \frac{dR}{dy},
\]
\[
\frac{d^{2}z}{dt^{2}} = - \frac{(M + M')\,z}{r^{3}} - \frac{dR}{dz},
\]
and the motion may be represented by supposing that the body moves in an ellipse
with variable elements, and such that the direction and velocity of motion are always
the same as in the actual orbit.

\newpage

%===============================================================
% PDF p.99 = book p.77
%===============================================================
\begin{center}
\small 176]\hfill\textsc{note on jacobi's canonical formul\ae.}\hfill\textsc{77}
\end{center}

\medskip

\noindent We may, to fix the ideas, take the plane of $xy$ to be the plane of the ecliptic
and the axis of $x$ to be the line through the first point of Aries, or origin of
longitude.

Jacobi's canonical elements may be taken to be,

First, the constant of vis viva or invariable part of the half-square of the
velocity, which is equal to the sum of the masses divided by twice the mean distance
and with the sign minus.

Secondly, the constant of areas, which is equal to the square root of the sum
of the masses into the half of the latus rectum.

Thirdly, the constant of the reduced area (i.e.\ of the area described on the plane
of the ecliptic), which is equal to the square root of the sum of the masses into the
half of the latus rectum into the cosine of the inclination.

Fourthly, the constant attached by addition to the time $t$, or what is the same
thing, the epoch or time of pericentric passage, taken with the sign minus.

Fifthly, the angular distance from node (or argument of latitude) of the pericentre.

Sixthly, the longitude of the node.

(The first and fourth elements are taken by Jacobi with the contrary sign, but
this difference is not material.)

Representing the preceding system of canonical elements by $\mathfrak{A}$, $\mathfrak{B}$, $\mathfrak{C}$, $\mathfrak{F}$, $\mathfrak{G}$, $\mathfrak{H}$,
and observing that Jacobi's disturbing function $\Omega$ is the same as the disturbing
function $R$ of the M\'ecanique C\'eleste, except that the sign is reversed (i.e.\ $\Omega = - R$),
the expressions for the variations of the canonical elements are
\[
\frac{d\mathfrak{A}}{dt} = - \frac{dR}{d\mathfrak{F}},\quad \frac{d\mathfrak{B}}{dt} = - \frac{dR}{d\mathfrak{G}},\quad \frac{d\mathfrak{C}}{dt} = - \frac{dR}{d\mathfrak{H}},
\]
\[
\frac{d\mathfrak{F}}{dt} = + \frac{dR}{d\mathfrak{A}},\quad \frac{d\mathfrak{G}}{dt} = + \frac{dR}{d\mathfrak{B}},\quad \frac{d\mathfrak{H}}{dt} = + \frac{dR}{d\mathfrak{C}}.
\]

In the ordinary case in which the disturbing force is the attraction of a third
body, write
\begin{align*}
& x',\ y',\ z',\quad\text{the coordinates of the disturbing body,}\\
& r',\quad\text{the radius vector, } = \sqrt{(x'^{2} + y'^{2} + z'^{2})},\\
& M'',\quad\text{the mass of the disturbing body.}
\end{align*}
Then the expression for the disturbing function is
\[
R = M''\left\{\frac{xx' + yy' + zz'}{r'^{3}} - \frac{1}{\sqrt{(x - x')^{2} + (y - y')^{2} + (z - z')^{2}}}\right\}.
\]

The preceding formulae form a convenient standard of reference for the various
systems of elements which have been made use of by writers upon Physical Astronomy;
any such system may be without difficulty derived from the canonical system by
expressing the elements adopted in terms of the above-mentioned canonical elements.

\newpage

%===============================================================
% PDF p.100 = book p.78 -- Paper 177 "Solution of a Mechanical Problem"
%===============================================================
\begin{center}
\small\textsc{78}\hfill\hfill[177
\end{center}

\vspace{4em}

\begin{center}
{\large\textbf{177.}}\\[1em]
{\large\scshape Solution of a Mechanical Problem.}\\[1em]
{\small [From the \textit{Quarterly Mathematical Journal}, vol.~\textsc{i}.\ (1857), pp.~405--406.]}
\end{center}

\bigskip

\textsc{A heavy} plane is supported by parallel elastic strings of small extensibility;
and the strings are of the same length and extensibility: required the position of
equilibrium.

Imagine the plane horizontal, and let $n$ be the number of strings, $(a, b)$, $(a', b')$,
$\&c.$ the coordinates of the points of attachment; $\xi$, $\eta$ the coordinates of the centre
of gravity of the plane; $W$ the weight; let the equation of the horizontal line about
which the plane turns be
\[
x\cos\alpha + y\sin\alpha - p = 0;
\]
and let $\delta\theta$ be the inclination of the plane in its position of equilibrium to the hori\-zontal plane, and $\omega\delta l$ the force generated by an increase $\delta l$ in the length of one of
the strings.

We have for the conditions of equilibrium
\[
\Sigma(a\cos\alpha + b\sin\alpha - p)\,\omega\delta\theta - W = 0,
\]
\[
\Sigma(a\cos\alpha + b\sin\alpha - p)\,a\omega\delta\theta - W\xi = 0,
\]
\[
\Sigma(a\cos\alpha + b\sin\alpha - p)\,b\omega\delta\theta - W\eta = 0;
\]
or putting $\Sigma a = L$, $\Sigma b = M$, $\Sigma a^{2} = A$, $\Sigma ab = H$, $\Sigma b^{2} = B$, we have
\[
L\cos\alpha + M\sin\alpha - np - \frac{W}{\omega\delta\theta} = 0,
\]
\[
A\cos\alpha + H\sin\alpha - Lp - \frac{W\xi}{\omega\delta\theta} = 0,
\]
\[
H\cos\alpha + B\sin\alpha - Mp - \frac{W\eta}{\omega\delta\theta} = 0.
\]

\end{document}
